\section{Введение}

На этой лекции мы продолжим работать со взвешенными структурами данных.
На ней мы решим задачу нахождения \textit{минимальных остовных деревьев}, а также смежную структуру данных -- \textit{систему непересекающихся множеств}, которая нам позволит быстрее выполнять поиск искомого дерева.

Всюду далее у нас есть граф $G\langle V, E \rangle$ и на множестве рёбер задана весовая функция $w: E \to \R$.

\section{Система непересекающихся множеств}

\Def Системой непересекающихся множеств назовём структуру данных, которая в любой момент поддерживает некоторый набор множеств $\mathcal{A} = \{A_1, A_2, \dots A_n\}: i \neq j \Rightarrow A_i \cap A_j = \varnothing$, а также позволяет выполнять три операции:
\begin{itemize}
    \item $\texttt{MAKE\_SET(}x\texttt{)}$ -- создает множество с единственным элементом $x$, который не принадлежит ни одному множеству.
    \item $\texttt{SAME(}x, y\texttt{)}$ -- возвращает \texttt{true}, если $\exists i: x, y \in A_i$ (то есть два элемента лежат в одном и том же множестве).
    \item $\texttt{UNITE(}x, y\texttt{)}$ -- выполняет объединение двух множеств, в которых лежат $x$ и $y$.
\end{itemize}
Элементы множества $A_i$ иногда назвают \textit{представителями множества}.

\Exe Предложите алгоритм, который выделяет компоненты связности, с помощью СНМ.

\subsection{Лес непересекающихся множеств}

Далее рассмотрим одну из реализаций СНМ -- \textit{лес непересекающихся множеств}.

Каждое множество будем представлять в виде дерева, которое состоит из элементов этого множества.

\begin{figure}[H]
    \centering
    \begin{tikzpicture}[
        node distance=1cm,
        state/.style={circle,draw,fill=white!20,minimum size=0.5cm}
    ]

    \node[state] (c) {$c$};
    \node[state] (h) [below left=of c] {$h$};
    \node[state] (e) [below right=of c] {$e$};
    \node[state] (b) [below=of h] {$b$};

    \draw[->,thick] (c) edge[loop above] (c);
    \draw[->,thick] (h) -- (c);
    \draw[->,thick] (e) -- (c);
    \draw[->,thick] (b) -- (h);


    \node[state] (f) [right=4cm of c] {$f$};
    \node[state] (d) [below=of f] {$d$};
    \node[state] (g) [below=of d] {$g$};

    \draw[->,thick] (f) edge[loop above] (f);
    \draw[->,thick] (d) -- (f);
    \draw[->,thick] (g) -- (d);

    \end{tikzpicture}
    \caption{Пример СНМ для $\mathcal{A} = \{ \{c, h, b, e\}, \{f, d, g\}\}$}
\end{figure}

Тогда в такой парадигме интерфейс реализуется следующим образом:
\begin{itemize}
    \item $\texttt{MAKE\_SET(}x\texttt{)}$ -- создание узла.
    \item $\texttt{SAME(}x, y\texttt{)}$ -- для $x$ и $y$ находим корни деревьев, в которых они лежат, $r_x$ и $r_y$ при помощи поднятия вверх. Затем сравниваем эти корни, если $x$ и $y$ лежат в одном множестве, то $r_x = r_y$, иначе $r_x \neq r_y$.
    \item $\texttt{UNITE(}x, y\texttt{)}$ -- также как в \texttt{SAME} находим корни $r_x$ и $r_y$, а затем подвешиваем $r_y$ в качестве ребенка к $r_x$.
\end{itemize}

Итак, исследуем время работы нашей структуры данных. Очевидно, \texttt{MAKE\_SET} работает за $\O(1)$.
Но вот с \texttt{UNIT} и \texttt{SAME} уже не всё так просто, потому что время их работы зависит напрямую от высот деревьев, в которых лежат $x$ и $y$.
И не сложно заметить, если взять два бамбука высоты $h$ и вызвать \texttt{SAME} от самых глубоких вершин, то время работы будет $\Theta(h)$.

На самом деле можно ускорить работу этих операций вплоть до <<почти константы>>. Для этого надо использовать две эвристики: \textit{ранговую} и \textit{сжатие путей}.

\subsubsection{Весовая эвристика}

Кажется достаточно разумным, чтобы уменьшить итоговую высоту поднятия, нужно подвешивать множество меньшего размера к большему.
Такая эвристика называется \textit{весовой}. Нам необходимо поддерживать размер текущего множества, что достаточно тривиально:
\begin{enumerate}
    \item При \texttt{MAKE\_SET} будет создавать новый узел и в счётчик ставить 1.
    \item При \texttt{UNITE} будем просто прибавлять к счётчику размера большего дерева счётчик меньшего.
\end{enumerate}

\Def Пусть $x$ -- элемент некоторого множества в СНМ. Тогда $s(x)$, по определению, положим размером данного множества, в котором лежит $x$.

\Def Пусть $x$ и $y$ -- элементы некоторого множества в СНМ. Тогда $x \cup y$, по определению, результат объединения множеств, в которых они лежат. $s(x \cup y)$ -- размер результрующего множества.

\Exe Пусть даны два множества в СНМ, первое больше второго. Верно ли утверждение, что высота первого больше высоты второго?

\Th Пусть выполнялось $m$ операций $\texttt{MAKE\_SET}, \texttt{SAME}, \texttt{UNITE}$, из которых $n$ операций $\texttt{MAKE\_SET}$. Тогда суммарное время работы при использовании весовой эвристики равно $\O(m \log n)$.
\begin{proof}
    Вспомним, что \texttt{MAKE\_SET} работает за $\O(1)$, поэтому его суммарное время выполнения есть $\O(n)$. 
    В тоже время время основной вклад в асимтотику дают \texttt{SAME}, \texttt{UNITE}, т.к. они напрямую зависят от того на какой глубине лежат рассматриваемые элементы $x$ и $y$, ведь мы <<прыгаем>> от них к корню.
    Можно сказать, что каждая из них работает за $\Theta(h(x) + h(y))$, где $h(v)$ -- глубина, на которой лежит вершина $v$. Далее приведём оценку $h$.

    Давайте рассмотрим произвольный элемент $x$ и исследуем его положение в течение всей последовательности операций. Изначально, когда только вызвали \texttt{MAKE\_SET}, он находился на глубине $h(x) = 0$.
    Далее рассмотрим случаи объединения с некоторым множеством, которое содержит некоторый элемент $y$:
    \begin{enumerate}
        \item $s(x) > s(y) \Rightarrow$ Тогда мы подвешиваем множество с элементом $y$ к множеству с элементом $x$ и $h(x)$ \textit{не меняется}. 
        \item $s(x) < s(y) \Rightarrow$ Тогда мы подвешиваем множество с элементом $x$ к множеству с элементом $y$. Но в этом случае глубина $x$ увеличилась на единицу. Также заметим, что $s(x \cup y) > 2 s(x)$.
        Действительно, т.к. мы подвесили множество с элементом $x$ к множеству с элементом $y$, которое не менее чем первое, то мы увеличились минимум в два раза. Более формально:
        \begin{equation*}
            s(x \cup y) = s(x) + s(y) > s(x) + s(x) = 2s(x)
        \end{equation*}
        \item $s(x) = s(y) \Rightarrow$ Предположим худший случай и мы подвесили множество с $x$ к множеству с $y$ (хотя можно было и наоборот, но опять же смотрим худший случай). Тогда высота $x$ опять увеличилась на единицу и размер результирующего увеличился в два раза:
        \begin{equation*}
            s(x \cup y) = s(x) + s(y) = s(x) + s(x) = 2s(x)
        \end{equation*}
    \end{enumerate}

    Из вышеописанного следует, что высота $x$ увеличивается в худшем случае на единицу. И она увеличивается, тогда и только тогда, когда это приводит к увелечению множества, в котором лежит $x$, минимум в два раза.
    Значит высота $x$ оценивается сверху количеством увечений его множества в два раза. А в наибольшем множестве, котором он мог находиться -- это в множестве размера $n$.
    Количество увелечений минимум в два раза до размера $n$ есть $\O(\log n)$. Значит, высота не более чем логарифмична.

    В частности из этого следует, что время работы \texttt{SAME} и \texttt{UNITE} есть $\O(\log n)$. Следовательно, в силу того, что у нас всего $m$ операций, то итоговая асимпотика есть $\O(m \log n)$.
\end{proof}

\Cons Пусть выполнялось $m$ операций $\texttt{MAKE\_SET}, \texttt{SAME}, \texttt{UNITE}$, из которых $n$ операций $\texttt{MAKE\_SET}$. Тогда амортизированная стоимость операций при использовании весовой есть $\O^*(\log n)$.
\begin{proof}

    \begin{equation*}
        \frac{\O(m \log n)}{m} = \O \left(\frac{m \log n}{m}\right) = \O^*(\log n)
    \end{equation*}
\end{proof}

\subsubsection{Ранговая эвристика\optstar}

Есть еще один способ аналогичный прошлому, его иногда используют как альтернативу. Чтобы реже приходить к ситуации бамбука и уменьшать высоту, надо подвешивать дерево с меньшей высотой к большему, потому что тогда высота почти не меняется.
Такой подход называется \textit{ранговой эвристикой}. Будем поддерживать ранг -- грубую оценку на высоту дерева. Ранг множества, в котором лежит вершина $x$, будем обозначать $r(x)$. Считается она следующим образом:
\begin{enumerate}
    \item При $\texttt{MAKE\_SET(}x\texttt{)}$ выставляем множеству ранг 0.
    \item При $\texttt{UNITE(}x, y\texttt{)}$ рассмотрим случаи:
        \begin{enumerate}
            \item $r(x) > r(y) \Rightarrow$ подвешиваем дерево, в котором лежит $y$, к дереву в котором лежит $x$. Ранг после этого равен $r(x)$ (ведь <<высота дерева>> не поменялась).
            \item $r(x) < r(y) \Rightarrow$ подвешиваем дерево, в котором лежит $x$, к дереву в котором лежит $y$. Ранг после этого равен $r(y)$ (ведь <<высота дерева>> не поменялась).
            \item $r(x) = r(y)$. Тут без разницы к кому что подвешивать, ведь ранг одинаков. Но после подвешивания ранг новоиспеченного множества будет равен $r(x) + 1$ (<<высота дерева>> увеличилась на единицу).
        \end{enumerate}
\end{enumerate}

\Exe Верно ли что если ранг одного дерева больше другого, то тоже верно для размера соответующих множеств?

\Exe Покажите, что если использовать только ранговую эвристику, то ранг есть максимальная глубина дерева.

Сама по себе эта эвристика дает уже прирост к работе. Следующие два утверждения даны для интересующегося темой читателя и \textit{не будут спрашиваться не экзамене}.

\ThBd Пусть выполнялось $m$ операций $\texttt{MAKE\_SET}, \texttt{SAME}, \texttt{UNITE}$, из которых $n$ операций $\texttt{MAKE\_SET}$. Тогда суммарное время работы при использовании ранговой эвристики равно $\O(m \log n)$.
% \begin{proof}

%     Вспомним, что \texttt{MAKE\_SET} работает за $\O(1)$, поэтому его суммарное время выполнения есть $\O(n)$. 
%     В тоже время время основной вклад в асимтотику дают \texttt{SAME}, \texttt{UNITE}, т.к. они напрямую зависят от того на какой глубине лежат рассматриваемые элементы $x$ и $y$, ведь мы <<прыгаем>> от них к корню.
%     Можно сказать, что каждая из них работает за $\Theta(h(x) + h(y))$, где $h(v)$ -- глубина, на которой лежит вершина $v$. Далее приведём оценку $h$.

%     Давайте рассмотрим произвольный элемент $x$ и исследуем его положение в течение всей последовательности операций. Изначально, когда только вызвали \texttt{MAKE\_SET}, он находился на глубине $h(x) = 0$.
%     Далее рассмотрим случаи объединения с некоторым множеством, которое содержит некоторый элемент $y$:
%     \begin{enumerate}
%         \item $r(x) > r(y) \Rightarrow$ Тогда мы подвешиваем множество с элементом $y$ к множеству с элементом $x \Rightarrow$ 
%     \end{enumerate}
    
% \end{proof}

\ConsBd Пусть выполнялось $m$ операций $\texttt{MAKE\_SET}, \texttt{SAME}, \texttt{UNITE}$, из которых $n$ операций $\texttt{MAKE\_SET}$. Тогда амортизированная стоимость операций $\texttt{MAKE\_SET}, \texttt{SAME}, \texttt{UNITE}$ есть $\O^*(\log n)$.
\begin{proof}

    \begin{equation*}
        \frac{\O(m \log n)}{m} = \O \left(\frac{m \log n}{m}\right) = \O^*(\log n)
    \end{equation*}

\end{proof}

То есть получается, уже сама по себе ранговая эвристика ускоряет нас с линейного до в каком-то плане логарифмического времени работы.

\subsubsection{Эвристика сжатия путей}

Также, кажется что неплохой идей будет узлы, для которых мы узнали корень (то есть все на пути при поднятии к корню), \textit{переподвешивать сразу к корню}.

Рассмотрим на примере работу данной эвристики. Пусть даны два дерева.

\begin{figure}[H]
    \centering
    \begin{tikzpicture}[
        node distance=1cm,
        state/.style={circle,draw,fill=white!20,minimum size=0.5cm,thick,font=\bfseries}
    ]

    \node[state] (b) {$b$};
    \node[state] (a) [below left=of b] {$a$};
    \node[state] (d) [below right=of b] {$d$};
    \node[state] (c) [below left=of d] {$c$};
    \node[state] (e) [below right=of d] {$e$};

    \draw[->,thick] (a) -- (b);
    \draw[->,thick] (d) -- (b);
    \draw[->,thick] (c) -- (d);
    \draw[->,thick] (e) -- (d);

    \begin{scope}[xshift=6cm]
        \node[state] (g) {$g$};
        \node[state] (f) [below left=of g] {$f$};
        \node[state] (h) [below right=of g] {$h$};
        \node[state] (i) [below right=of h] {$i$};
        
        \draw[->,thick] (f) -- (g);
        \draw[->,thick] (h) -- (g);
        \draw[->,thick] (i) -- (h);
    \end{scope}

    \end{tikzpicture}
\end{figure}

И мы хотим проверить что $e$ и $i$ лежат в одном множестве. Тогда после вызова $\texttt{SAME(}e, i\texttt{)}$ конфигурация деревьев будет следующая:

\begin{figure}[H]
    \centering
    \begin{tikzpicture}[
        node distance=1cm,
        state/.style={circle,draw,fill=white!20,minimum size=0.5cm,thick,font=\bfseries}
    ]

    \node[state] (b) {$b$};
    \node[state] (a) [below left=of b] {$a$};
    \node[state] (d) [below right=of b] {$d$};
    \node[state] (e) [right=of d] {$e$};
    \node[state] (c) [below left=of d] {$c$};

    \draw[->,thick] (a) -- (b);
    \draw[->,thick] (d) -- (b);
    \draw[->,thick] (e) -- (b);
    \draw[->,thick] (c) -- (d);

    \begin{scope}[xshift=7cm]
        \node[state] (g) {$g$};
        \node[state] (f) [below left=of g] {$f$};
        \node[state] (h) [below right=of g] {$h$};
        \node[state] (i) [right=of h] {$i$};
        
        \draw[->,thick] (f) -- (g);
        \draw[->,thick] (h) -- (g);
        \draw[->,thick] (i) -- (g);
    \end{scope}

    \end{tikzpicture}
\end{figure}

Тогда следующий вызов $\texttt{SAME(}e, i\texttt{)}$, очевидно, будет быстрее, ведь нам нужно подняться на один уровень вместо 2-х. 
Также как и ранговая, сама по себе эвристика сжатия путей тоже дает выигрыш. Также как и с ранговой эвристикой, следующее утверждение на экзамене спрашиваться не будет.

\ThBd Если было выполнено $n$ операций $\texttt{MAKE\_SET}$, а также $m$ операций \texttt{SAME}. Тогда суммарное время работы при использовании ранговой эвристики равно $\O\left(n + \frac{m}{2} \log_{2 + \frac{m}{2n}} n \right)$.

\subsubsection{Объединение по рангу со сжатием путей}

Самый большой прирост даёт испольвание этих двух эвристик одновременно. Если раньше мы могли амортизированно дойти до логарифма, то с двумя эвристиками мы можем дойти до почти константы.
Давайте же рассмотрим эту <<почти константу>>.

\Def Функцией Аккермана назовём следующую функцию:
\begin{equation*}
    A(m, n) \overset{def}{=} \begin{cases}
        n + 1, m = 0 \\
        A(m - 1, 1), m > 0, n = 0\\
        A(m - 1, A(m, n - 1)), m > 0, n > 0
    \end{cases}
\end{equation*}

Функция $A$ является чрезывчайно быстро растущей функцией.

\Def Обратной функцией Аккермана положим $\alpha(n) \overset{def}{=} \min\limits_{k \in \N} \{ k: A(k, k) \geq n  \}$.

На практике $\alpha(n) \leq 4$, поэтому в рамках разумных задач, можно считать её константой (но теоретически, все же это функция).

Следующая теорема даёт нам желанный результат. Сие доказательство столь сложно, что Кормен расположил его на 10 страницах А4, поэтому, очевидно, мы принимаем его без доказательства.

\ThBd Пусть выполнялось $m$ операций $\texttt{MAKE\_SET}, \texttt{SAME}, \texttt{UNITE}$, из которых $n$ операций $\texttt{MAKE\_SET}$. 
Тогда при использовании обоих эвристик время выполнения есть $\O(m \alpha(n))$.

\ConsBd Пусть выполнялось $m$ операций $\texttt{MAKE\_SET}, \texttt{SAME}, \texttt{UNITE}$, из которых $n$ операций $\texttt{MAKE\_SET}$. Тогда амортизированная стоимость операций есть $\O^*(\alpha(n))$.
\begin{proof}
    \begin{equation*}
        \frac{\O(m \alpha(n))}{m} = \O\left(\frac{m \alpha(n)}{m}\right) = \O^*\left(\alpha(n)\right)
    \end{equation*}
\end{proof}

\section{Минимальные остовные деревья}

\Def Пусть $G \langle V, E \rangle$ -- граф. Тогда граф $G' \langle V', E' \rangle$, будем называть \textit{подграфом}, если $V' \subset V, E' \subset E$. Если граф $G'$ является подграфом $G$, будем обозначать это $G' \subset G$.

\Def Пусть $T \subset G$. Если $T$ -- дерево, тогда $G'$ -- \textit{остовное дерево} графа $G$. Остовное дерево также называют остовом.

\Def Пусть $G \langle V, E \rangle$ -- граф, $w: E \to \R$ -- весовая функция. Тогда \textit{весом графа} назовём следующую величину:
\begin{equation*}
    w(G) \overset{def}{=} \sum\limits_{e \in E} w(e)
\end{equation*}

\Def Пусть $G \langle V, E \rangle$ -- связный граф, $w: E \to \R$ -- весовая функция. \textit{Минимальным остовным деревом} назовём такое остовное дерево, вес которого минимален среди всех друг остовных деревьев графа $G$. Иногда минимальное остовное дерево называют \textit{миностовом}.

Собственно задача, которую мы будем решать оставшуюся часть лекции -- это алгоритмы построения тех самых миностовов. 
Все алгоритмы, которые мы разберём, по своей сути жадные и используют под капотом одну и ту же лемму, которую мы прямо сейчас докажем.

Далее под графом мы везде позразумеваем связный граф, если не оговорено обратное (ведь остовные деревья не в связных графах не имеют смысла).

\subsection{Лемма о безопасном ребре}

\Note Далее мы будем отождествлять множество рёбер $E$ с соответствующим графом $G \langle V, E \rangle$, множество вершин которого совпадает с множеством вершин, используемых в рёбрах из $E$.

Всюду далее мы считаем, что у нас задан некоторый граф $G \langle V, E \rangle$ и весовая функция $w: E \to \R$.

\Def Пусть $A$ -- некоторое множество рёбер, которое является подграфом некоторого миностова. Пусть $e \in E$. Ребро $e$ будем называть \textit{безопасным} относительно $A$, если $A \cup \{e\}$ -- подграф некоторого миностова.

Ясно, что такие ребра чрезвычайно полезны и из определения безопасного ребра уже в принципе можно сделать алгоритмы, которые строят миностов. 
Но все это упирается в то, как искать вообще такие рёбра. Следующая лемма немного нам прояснит как они ведут себя.

\Def Разрезом будем называть пару $(S, T): S \cap T = \varnothing, S \cup T = V$.

\Def Будем говорить, что ребро $(u, v) \in E$ пересекает разрез $(S, T)$, если $u \in S$ и $v \in T$, либо $u \in T$ и $v \in S$.

\Lm{О безопасном ребре} Пусть множество рёбер $A$ -- подграф некоторого миностова. Пусть $(S, T)$ -- разрез $G$ таков, что ни одно ребро из $A$ не пересекает разрез $(S, T)$.
Если $e \in E$ -- ребро минимального веса из тех, что пересекают разрез $(S, T)$, то $e$ -- безопасное ребро для $A$.
\begin{proof}
    
    Пусть $T$ -- некоторый миностов, что $A \subset T$ (такой существует, по определению $A$). Если $A \cup \{e\} \subset T$, то теорема доказана.
    Рассмотрим случай, когда $A \cup \{e\} \not\subset T$.  Будем считать, что $e = (u, v)$ и $u \in S, v \in T$.

    Заметим, что при добавлении $e$ в $T$ образуется цикл (пример этого см. на рисунке ниже). Также заметим, что  в этом цикле есть хотя бы одно ребро пересекающее разрез $(S, T)$ отличное от $(u, v)$ (на рисунке это ребро $(b, c)$).
    Обозначим данное ребро за $e'$.

    \begin{figure}
        \centering
       \begin{tikzpicture}[
            main_node/.style={circle, draw, thick, minimum size=6mm, inner sep=0pt, fill=white},
            edge_style/.style={thick},
            oval_style/.style={ellipse, draw, very thick, inner xsep=1.5cm, inner ysep=1.5cm},
            blue_line_style/.style={blue, very thick},
            highlight_edge/.style={red, line width=2mm, opacity=0.2, line cap=butt},
            label_style/.style={font=\Large\bfseries}
        ]

        \coordinate (center_left) at (-2.5, 0);
        \coordinate (center_right) at (2.5, 0);

        \node[oval_style] (oval_L) at (center_left) {};
        \node[oval_style] (oval_R) at (center_right) {};

        \node[label_style, below=0.5cm of oval_L.south] {$S$};
        \node[label_style, below=0.5cm of oval_R.south] {$T$};


        \node[main_node] (a) at ($(center_left) + (-0.5, 1)$) {$a$};
        \node[main_node] (b) at ($(center_left) + (-0.5, -0.5)$) {$b$};
        \node[main_node] (u) at ($(center_left) + (1.5, 0.5)$) {$u$};

        \node[main_node] (v) at ($(center_right) + (-1.5, 0.5)$) {$v$};
        \node[main_node] (c) at ($(center_right) + (1.0, -0.5)$) {$c$};

        \draw[edge_style] (a) -- (u);
        \draw[edge_style] (b) -- (u);
        \draw[edge_style] (b) -- (c);
        \draw[edge_style] (v) -- (c);

        \draw[blue_line_style] (u) -- (v);

        \draw[highlight_edge] ($(v)!3mm!(c)$) -- ($(c)!3mm!(v)$);
        \draw[highlight_edge] ($(c)!3mm!(b)$) -- ($(b)!3mm!(c)$);
        \draw[highlight_edge] ($(b)!3mm!(u)$) -- ($(u)!3mm!(b)$);

        \end{tikzpicture} 

        \caption{Пример цикла, который образовался при добавлении ребра $(u, v)$}
    \end{figure}

    Удалим ребро $e'$ из $T$. Как следствие $T$ разделяется на две компоненты связности. Теперь возьмем и соединим эти две компоненты ребром $e$, получив некоторое новое остовное дерево $T'$ (на примере ниже, удаляем $(b, c)$ и вставляем $(u, v)$).

    Далее, заметим, что:
    \begin{equation*}
        w(T') = w(T) - w(e') + w(e)
    \end{equation*}

    Также $e'$ -- ребро пересекающее разрез при том, что $e$ выбрано как ребро минимального веса из всех пересекающих разрез. Из этого следует, что:
    \begin{equation*}
        w(e) - w(e') \leq 0
    \end{equation*}
    Комбинируя с предыдущим равенством, получаем:
    \begin{equation*}
        w(T') = w(T) - w(e') + w(e) \leq w(T)
    \end{equation*}
    \begin{equation*}
        w(T') \leq w(T)
    \end{equation*}
    Но в силу того, что $T$ -- остовное дерево минимального веса, то $T'$ -- тоже миностов. Теперь покажем, что $A \cup \{e\}$ есть подграф $T'$.
    Заметим, что $e' \not\in A$, по определению разреза $(S, T)$. Поэтому если отождествить $T$ и $T'$ с множеством соответствующих рёбер, то получаем следующие соотношения:
    \begin{equation*}
        A \subset T \setminus \{ e' \} \Rightarrow A \cup \{ e \} \subset (T \setminus \{e'\}) \cup \{ e \} = T'
    \end{equation*}
    \begin{equation*}
        A \cup \{ e \} \subset T'
    \end{equation*}

    Получили, что $A \cup \{e\}$ есть подграф некоторого миностова $T'$, значит, по определению, $e$ -- безопасное.

\end{proof}
